\section{生产环境配置规范与 YAML 范例}
\label{app:config}

\textbf{Private LLM Analytics} 系统遵循声明式配置与代码逻辑严格解耦的软件工程原则。所有隐私界限、机制参数、记账策略、外部接收方权限与领域分类均通过标准 YAML 文件进行配置驱动。本附录提供与当前仓库 YAML 对齐的核心配置节选；不是对所有入口默认值已统一的声明。

\subsection{指标机制定义配置 (metrics.yaml)}

\texttt{configs/metrics.yaml} 统一定义系统中所有已注册统计指标的维度结构、用户级贡献限制模式与噪声释放机制类型。

\begin{techbox}{指标机制生产配置清单}
\begin{verbatim}
metrics:
  # 1. 用户使用率指标
  user_prevalence:
    enabled: true
    type: prevalence
    dimensions:
      - provider
      - task
    contribution:
      mode: binary              # 二值指示向量
      max_cells_per_user: 9     # 单用户最多参与单元数截断界限 k
    mechanism:
      type: gaussian

  # 2. 用户归一化使用占比指标
  user_normalized_share:
    enabled: true
    type: normalized_share
    dimensions:
      - provider
      - task
    contribution:
      mode: l1_normalized       # L1 归一化用户工作流向量
      norm_bound: 1.0           # 天然敏感度上界
    mechanism:
      type: gaussian

  # 3. 交互调用总量指标
  interaction_volume:
    enabled: true
    type: volume
    dimensions:
      - provider
      - task
    contribution:
      mode: l2_clip             # 用户多维调用频次 L2 范数截断
      clip_norm: 20.0           # 截断界限 C_vol
    mechanism:
      type: gaussian

  # 4. 任务成功率指标
  success_rate:
    enabled: true
    type: success_rate
    dimensions:
      - provider
      - task
    contribution:
      mode: joint_l2_clip       # 共享调用总量总范数的联合截断
      clip_norm_rate: 20.0      # 成功率截断界限 C_rate
      min_noisy_denominator: 1.0 # 有噪分母抑制最小阈值 tau
    mechanism:
      type: gaussian

  # 5. 热门任务发现指标
  top_k_tasks:
    enabled: true
    type: top_k_tasks
    dimensions:
      - task
    contribution:
      mode: binary
      max_tasks_per_user: 5
      top_k: 5                  # T=3 时实际最多输出 min(K,T)=3 项
    mechanism:
      type: gaussian
\end{verbatim}
\end{techbox}

本文所有 Success Rate 配置示例均显式设置 \texttt{min\_noisy\_denominator: 1.0}。YAML 加载后和 Profile 默认值为 $1.0$；部分兼容构造入口仍默认 $5.0$，因此不可省略此字段。Top-K 的配置 $K=5$ 保留当前代码行为，实际输出数受任务字典大小限制。

\subsection{全生命周期隐私与规划器配置 (privacy.yaml)}

\texttt{configs/privacy.yaml} 集中配置全系统的差分隐私基础假设、生命周期预算上限、RDP 记账阶数、外部接收方合并模式以及数据驱动离线规划器（Offline Planner）的先验分布。

\begin{techbox}{隐私管理与记账生产配置清单}
\begin{verbatim}
privacy:
  privacy_unit: user            # 隐私保护实体：单个自然人用户
  adjacency: add_remove         # 邻接关系：加/删单用户所有记录

  horizon:
    days: 90                    # 规划与记账生命周期时间跨度

  global_budget:
    epsilon: 4.0                # basic 模式/内部作用域的 epsilon 上限
    delta: 4.0e-6               # basic 模式/内部作用域的累加 delta 上限

  default_query_budget:
    epsilon: 1.0                # 单次常规发布默认隐私预算
    delta: 1.0e-6

  accountant:
    mode: rdp                   # 采用 Rényi 差分隐私紧致组合记账
    epsilon_cap: 4.0
    delta_global: 1.0e-6         # 外部 RDP 转换和 lifetime cap 的 delta
    # 跟踪的 Rényi 散度阶数列表 alpha
    orders:
      - 1.25
      - 1.5
      - 2.0
      - 3.0
      - 4.0
      - 5.0
      - 8.0
      - 10.0
      - 16.0
      - 32.0
      - 64.0
      - 128.0

  history:
    mode: lifetime              # 永久追踪全部发布历史事件

  recipient_composition:
    mode: global_external       # 所有外部模型提供商共用同一外部预算池

  default_mechanism:
    type: gaussian              # 核心机制：解析高斯机制 (Balle-Wang)

  release:
    # 开启规范发布语义缓存，相同查询零预算开销
    reuse_existing_release: true
    allow_budget_overrun: false    # 预算耗尽时执行 Fail-Closed 熔断

  planner:
    enabled: true
    mode: static_public         # 基于公开先验参数的离线规划
    prior:
      distribution: contribution_model
      mean: 60.0
      std: 20.0
      lower: 5.0
      upper: 150.0
      public_seed: 900001
      num_reference_users: 2000
    weights:
      objective_mode: EXPECTED_RELEASE_L1
      data: 1.0
      clipping: 1.0
      noise: 1.0
      aggregation: 1.0
      staleness: 0.5
      window: 0.2
    candidate_clipping: [5.0, 10.0, 20.0, 40.0, 80.0]
    candidate_windows: [7, 14, 30, 60, 90]
    candidate_frequencies: [1, 7, 14, 30]
\end{verbatim}
\end{techbox}

\textbf{字段语义：}外部 \texttt{mode: rdp} 发布以 \texttt{accountant.delta\_global} 为唯一 RDP 转换目标，cap 为 \texttt{accountant.epsilon\_cap}。静态 \texttt{internal} 作用域在当前编排器中走 basic 组合，使用 \texttt{global\_budget} 的 $4.0,4\times10^{-6}$；显式 basic 模式亦使用这组值。两组 delta 含义不同，不能相加或互换。单次 \texttt{default\_query\_budget.delta} 只用于机制标定；持续观测使用 RDP 路径及其配置。

\subsection{多接收方权限与作用域配置 (recipients.yaml)}

\texttt{configs/recipients.yaml} 实现了基于最小权限原则的强隔离策略。系统将接收方划分为内部受信任审计方与外部竞争性模型服务商，并对外部实体施加列级过滤与跨供应商视图屏蔽。

\begin{techbox}{接收方角色与权限矩阵配置}
\begin{verbatim}
recipients:
  # 内部受信任治理团队
  internal:
    role: trusted
    display_name: "内部 AI 治理团队"
    providers: ["*"]            # 可查看所有模型供应商
    metrics: ["*"]              # 可访问全部指标体系
    exact_results: true         # 允许在受控沙盒内访问精确原始统计
    dp_results: true
    cross_provider_comparison: true

  # 外部模型供应商 A
  provider_a:
    role: external
    display_name: "OpenAI"
    providers:
      - provider_a              # 仅允许查看其自身模型对应的统计切片
    metrics:
      - user_prevalence
      - user_normalized_share
      - interaction_volume
      - success_rate
    exact_results: false        # 严格禁止访问明细或精确基准值
    dp_results: true            # 仅允许接收差分隐私脱敏释放结果
    cross_provider_comparison: false # 强隔离：屏蔽所有横向对比字段

  # 外部模型供应商 B
  provider_b:
    role: external
    display_name: "Kimi"
    providers:
      - provider_b
    metrics:
      - user_prevalence
      - user_normalized_share
      - interaction_volume
      - success_rate
    exact_results: false
    dp_results: true
    cross_provider_comparison: false

  # 外部模型供应商 C
  provider_c:
    role: external
    display_name: "DeepSeek"
    providers:
      - provider_c
    metrics:
      - user_prevalence
      - user_normalized_share
      - interaction_volume
      - success_rate
    exact_results: false
    dp_results: true
    cross_provider_comparison: false
\end{verbatim}
\end{techbox}

\subsection{分类体系与基准场景配置 (taxonomy.yaml \& scenario.yaml)}

系统通过独立分类配置实现对多模型及业务任务的动态扩展能力，避免在代码内部硬编码业务标签。

\begin{techbox}{分类与基准场景配置}
\begin{verbatim}
# configs/taxonomy.yaml
taxonomy:
  version: 1
  ordering: "canonical_lexicographic" # 强制采用规范字典序对齐多维向量
  providers:
    - provider_a
    - provider_b
    - provider_c
  tasks:
    - technical
    - document
    - analysis

# configs/scenario.yaml
dataset:
  source: synthetic
population:
  mode: fixed_public_population
  n_users: 50000              # 企业级合成基准总体规模：50,000 名合成用户
  n_days: 30                  # 观测周期跨度：30 天
usage_frequency:
  distribution: normal
  mean: 60                    # 单用户平均调用频次基准
  std: 20
  minimum: 5
  maximum: 150
random_seed: 42
\end{verbatim}
\end{techbox}
